\titledquestion{Array Representation of Min-Heap}
We use an array starting at index 0 to store a min complete binary heap, with the root of the heap placed at the index 1.
\begin{parts}


\part[2] Given an array $[0,1,2,8,5,30,10,12,60,7,81,78]$, which represents a min complete binary heap and the entry $1$ is the root at index $1$. Perform the \textbf{pop} operation according to the pop rules discussed in class. After the operation, the array is:
\par\noindent $[0, 2, 5, 8, 7, 30, 10, 12, 60, 78, 81]$

\part[2] Given an array $[0,1,2,8,5,30,10,12,60,7,81,78]$, which represents a min complete binary heap and the entry $1$ is the root at index $1$. Perform the \textbf{push 0} operation according to the pop rules discussed in class. After the operation, the array is:
\par\noindent $[0, 0, 2, 1, 5, 30, 8, 12, 60, 7, 81, 78, 10]$

\part Complete the following C++ skeleton of \textbf{Push} and \textbf{Pop}.
\begin{cpp}
void push(vector<int>& heap, int val) {
    heap.push_back(val); 
    int index = heap.size() - 1;  
    while (index > 1) {
        int parent = ____1____;  
        if (heap[parent] <= heap[index]) {
            break; 
        }
        _______2_______;
        _______3_______; 
    }
}

void pop(vector<int>& heap) {
    heap[1] = ____4____; 
    heap.pop_back(); 
    int index = 1;
    while (index < heap.size()) {
        int left = ____5____;  
        int right = ___6_____;  
        int smallest = index;
        if (left < heap.size() && heap[left] < heap[smallest]) {
            smallest = left;
        }
        if (right < heap.size() && heap[right] < heap[smallest]) {
            smallest = right;
        }
        if (smallest != index) {
            _______7_______;
            _______8_______; 
        } 
        else {
            break; 
        }
    }
}
\end{cpp}
\newpage
\begin{subparts}
    \subpart[2] Choose one of the following statements to fill in the blank 2,3,7,8:
    
        
    \vspace{5pt}
    \begin{checkboxes}
  \choice \begin{lstlisting}[language=C++]
  std::swap(heap[parent], heap[index]); //2
  index = heap.size()-1;   //3
  std::swap(heap[smallest], heap[index]); //7
  index = std::min(left,right); //8
  \end{lstlisting}
  
  \choice \begin{lstlisting}[language=C++]
  index = heap.size()-1;   //2
  std::swap(heap[parent], heap[index]); //3
  index = std::min(left,right); //7
  std::swap(heap[smallest], heap[index]); //8
  \end{lstlisting}

  \CorrectChoice \begin{lstlisting}[language=C++]
  std::swap(heap[parent], heap[index]); //2
  index = parent;   //3
  std::swap(heap[smallest], heap[index]); //7
  index = smallest; //8
  
  \end{lstlisting}
  
  \choice \begin{lstlisting}[language=C++]
  index = parent; //2
  std::swap(heap[parent], heap[index]); //3
  index = smallest; //7
  std::swap(heap[smallest], heap[index]); //8
  
  \end{lstlisting}
\end{checkboxes}

    
    \vspace{15pt}
    \setlength{\parskip}{0.6\baselineskip}
    \subpart[1] Fill in the blank (1):
    \par\noindent index / 2
    
    \vspace{15pt}
    \subpart[1] Fill in the blank (4):
    \par\noindent heap[heap.size() - 1]

    \vspace{15pt}
    \subpart[1] Fill in the blank (5):
    \par\noindent index * 2

    \vspace{15pt}
    \subpart[1] Fill in the blank (6):
    \par\noindent index * 2 + 1
\end{subparts}



\newpage
\part Complete the following C++ skeleton of building a heap through Floyd’s method.
\begin{cpp}
void heapify(vector<int>& heap, int i) {
    int n = heap.size() - 1;
    int left = ...   #dont need answer here
    int right = ...  #dont need answer here
    int smallest = i;        
    if (left <= n && heap[left] < heap[smallest]) {
        smallest = left;
    }
    if (right <= n && heap[right] < heap[smallest]) {
        smallest = right;
    }
    if (smallest != i) {
        _______1_______;
        _______2_______;  
    }
}
void buildHeap(vector<int>& heap) {
    int n = heap.size() - 1;
    for (int i = _______3_______) #Only necessary iterations, not 1...n
    {
        heapify(heap, i);  
    }
}
\end{cpp}
\begin{subparts}

    \vspace{15pt}
    \setlength{\parskip}{0.6\baselineskip}
    \subpart[1] Fill in the blank (1):
    \par\noindent std::swap(heap[i], heap[smallest])
    
    \vspace{15pt}
    \subpart[1] Fill in the blank (2):
    \par\noindent heapify(heap, smallest)

    \vspace{15pt}
    \subpart[1] Fill in the blank (3):
    \par\noindent n / 2; i >= 1; i--

\end{subparts}

\newpage
\part[3] Assume that all keys are distinct and that the input keys are generated in a uniformly random order. The heap is obtained by inserting the first $n$ keys in this order, and the new entry is the $(n+1)$-st key.

Write down the \textbf{average} time complexity of the following operations in a min complete binary heap of size $n$: (don't need to prove; use $\Theta$):
\begin{itemize}
    \item[1.] Access the root.
    \item[2.] Pop the root.
    \item[3.] Build a heap by Floyd’s Method.
    \item[4.] Merge two binary heaps of size $n$.
    \item[5.] Push a randomly generated entry.
    \item[6.] Push an entry smaller the root of the heap.
\end{itemize}
\begin{solution}
\begin{enumerate}
    \item[1.] $\Theta(1)$
    \item[2.] $\Theta(\log n)$
    \item[3.] $\Theta(n)$
    \item[4.] $\Theta(n)$
    \item[5.] $\Theta(1)$
    \item[6.] $\Theta(\log n)$
\end{enumerate}
\end{solution}
\end{parts}
